\documentclass[12pt,a4paper]{article}

\usepackage[utf8]{inputenc}
\usepackage[T1]{fontenc}
\usepackage[spanish,es-noshorthands,es-noindentfirst]{babel}
\usepackage{amsmath}
\usepackage{graphicx}
\usepackage{float}
\usepackage{enumitem}
\usepackage{listings}
\usepackage{xcolor}
\usepackage[margin=2.5cm]{geometry}
\usepackage[colorlinks=true,linkcolor=cyan!70!blue,urlcolor=cyan!70!blue]{hyperref}

\setlist[enumerate]{label=\arabic*.}

% --- Si una imagen no existe, se muestra un recuadro en su lugar ---
% Carpeta de imágenes (usar "/" y terminar con "/")
\newcommand{\rutaimg}{IMAGES/}
\newcommand{\imagen}[2][\linewidth]{%
  \IfFileExists{\rutaimg #2}{\includegraphics[width=#1]{\rutaimg #2}}%
  {\fbox{\parbox{#1}{\centering\vspace{2cm}Falta la imagen: \texttt{#2}\vspace{2cm}}}}%
}

% --- Estilo de listados de código ---
\lstset{
  basicstyle=\ttfamily\footnotesize,
  numbers=left,
  numberstyle=\tiny\color{gray},
  breaklines=true,
  frame=single,
  columns=fullflexible,
  keepspaces=true
}

\title{%
  \imagen[0.5\linewidth]{logo_unmsm.png}\\
  Laboratorio Dirigido Nro. 03 : IGBTs}
\author{Profesora: MsC. Luz Adanaqué}
\date{October 7, 2026}

\begin{document}

\maketitle

%--------------------------------------------------------------
\section{Indicaciones}
\begin{itemize}
  \item La entrega principal de este laboratorio consta de la tabla de
        resultados completa y las gráficas comparativas.
  \item Sea claro en sus simulaciones en LTspice; las señales de prueba
        ($V_{CE}$, $I_C$, $V_{GE}$) deben presentarse de manera ordenada,
        priorizando la visualización correcta de los transitorios.
  \item El día 21 se realizará una breve exposición con los principales
        hallazgos en su experiencia de laboratorio, en base a ello se
        ponderará la calificación.
  \item El laboratorio puede ser en parejas o individual.
\end{itemize}

%--------------------------------------------------------------
\section{Título}
Caracterización y Simulación Comparativa de Transistores IGBT en LTspice

%--------------------------------------------------------------
\section{Herramientas}
\begin{enumerate}
  \item LTspice.
  \item Hojas de datos y modelos SPICE.
  \begin{enumerate}[label=(\alph*)]
    \item Duncan Electronics Theorical Models: NIGBT y PIGBT.
    \item Fuji Electric Serie X: FGW50XS65, FGW50XS65C, FGW75XS65,
          FGW75XS65C.
  \end{enumerate}
\end{enumerate}

Esto se puede encontrar en los siguientes enlaces:
\begin{enumerate}[label=(\alph*)]
  \item \url{https://www.duncanamps.com/spicemosfet.html}
  \item \url{https://www.fujielectric.com/products/semiconductors_devices/\allowbreak
        semiconductor/product_series/\allowbreak
        model_power_discrete_technical_discrete_p-spice.html}
\end{enumerate}

%--------------------------------------------------------------
\section{Agregar Componentes de Modelos SPICE en LTspice}
Para agregar los modelos SPICE dentro del entorno de LTSpice, seguir el siguiente procedimiento.

\begin{enumerate}
  \item Guardar el archivo descargado en el mismo directorio donde se creará
        el esquema de LTspice (\texttt{.asc}).
  \item Añadir la directiva de texto \texttt{.include nombre\_del\_modelo.lib}
        directamente en el diseño esquemático (Como en la figura~\ref{fig:nigbt}).
  \item Insertar un símbolo genérico de IGBT desde la biblioteca de
        componentes.

  \begin{figure}[H]
    \centering
    \imagen[0.4\linewidth]{fig1_nigbt_lib.png}
    \caption{Ejemplo de NIGBT custom mediante librería.}
    \label{fig:nigbt}
  \end{figure}

  \item Hacer clic derecho sobre el símbolo del IGBT y modificar el atributo
        \texttt{Value} para que coincida exactamente con el nombre del modelo
        \texttt{.subckt} definido dentro del archivo de texto del fabricante.
        En la lista~\ref{lst:pigbt} se muestra el nombre del subcircuito
        resaltado para el modelo \texttt{PIGBT.lib}.
\end{enumerate}

\begin{lstlisting}[caption={Modelo SPICE del P-IGBT genérico},label={lst:pigbt},escapeinside={(*@}{@*)}]
*PIGBT
*P-IGBT Default MCE C G E 5-20-95
*MCE 600V 40A 24.3ns pkg:TO-220
.SUBCKT XPIGBT 71 72 74
Q1 83 81 85 QOUT
M1 81 82 83 83 MFIN L=1U W=1U
DSD 81 83 DO
DBE 81 85 DE
RC 85 71 21.2M
RE 83 73 2.12M
RG 72 82 25.5
CGE 82 83 1.42N
CGC 82 71 1P
EGD 91 0 81 82 1
VFB 93 0 0
FFB 81 82 VFB 1
CGD 92 93 1.41N
R1 92 0 1
D1 91 92 DLIM
DHV 94 93 DR
R2 91 94 1
D2 94 0 DLIM
LE 73 74 7.5N
DLV 94 95 DR 13
RLV 95 0 1
ESD 96 93 POLY(1) 81 83 19 1
MLV 95 96 93 93 SW
.MODEL SW NMOS (LEVEL=3 VTO=0 KP=5)
.MODEL QOUT NPN (IS=37.7F NF=1.2 BF=5.1 CJE=3.48N TF=24.3N XTB=1.3)
.MODEL MFIN PMOS (LEVEL=3 VMAX=400K THETA=46.1M ETA=2M VTO=-5.2 KP=2.12)
.MODEL DR D (IS=37.7F CJO=100P VJ=1 M=.82)
.MODEL DO D (IS=37.7F BV=600 CJO=2.07N VJ=1 M=.7)
.MODEL DE D (IS=37.7F BV=14.3 N=2)
.MODEL DLIM D (IS=100N)
.ENDS XPIGBT
\end{lstlisting}

\begin{enumerate}[resume]
  \item Si se presenta un error como el que se muestra en la
        lista~\ref{lst:error}, hacer \texttt{Ctrl + Click Derecho} en el
        componente, y en el apartado \texttt{Prefix} cambiar la \texttt{Z}
        por una \texttt{X}.
\end{enumerate}

\begin{lstlisting}[caption={Error de Prefijo},label={lst:error}]
LTspice 26.0.1 for Windows
Circuit: C:\Users\user\Downloads\NIGBT.net
C:\Users\user\Downloads\NIGBT.net (3): Undefined model "xirgbc20s".
Note that there is a sub-circuit with the same name "XIRGBC20S" defined,
which might indicate that this error is simply caused by the wrong prefix letter.
Consider using the prefix "X" instead of "Z".
You can change the prefix by editing the "Prefix" attribute in the component attribute editor.
Z1 N001 N004 0 XIRGBC20S
               ^^^^^^^^^
\end{lstlisting}

%--------------------------------------------------------------
\section{Procedimiento}

\subsection{Parte I: Simulación comparativa entre IGBT de canal N y canal P}
En esta sección se detallan los pasos para implementar el banco de pruebas
Double Pulse Test (DPT) para ambos tipos de transistores.

\begin{enumerate}
  \item \textbf{Configuración del circuito para el N-IGBT (Lado Bajo / Low-Side)}\\
        Realice el esquema de conexión mostrado en la figura~\ref{fig:esquema}.

        \textbf{Explicación:}
        \begin{enumerate}
          \item Terminal del \textbf{Emisor} conectada a la referencia (Tierra / 0\,V)
                mediante una resistencia de $R = 0.001\,\Omega$ (para medir
                corriente).
          \item Inductor de carga en serie con la fuente de voltaje conectado
                al \textbf{Colector}, esta fuente debe ser una fuente de corriente
                continua (VDC) y de un valor alto.
          \item Colocar un diodo de protección en paralelo con el inductor.
          \item Configure una fuente de voltaje de pulsos en la Compuerta
                ($V_{GE}$). Conectar en serie con una resistencia $R_G$ de
                ej.\ $10\,\Omega$.
        \end{enumerate}

  \item \textbf{Configuración del circuito para el P-IGBT (Lado Alto / High-Side)}\\
        Realice el esquema de conexión mostrado en la figura~\ref{fig:esquema}.
        \begin{enumerate}
          \item Inserte el símbolo del P-IGBT y asigne su modelo SPICE.
          \item Realice una configuración similar a la del N-IGBT.
          \item Configure la fuente de disparo $V_{GE}$ entre el terminal Gate
                (a través de $R_G$) y el \textbf{Emisor}. En modo \texttt{PULSE}, los niveles
                deben ser de 0\,V (apagado) a $-15$\,V (encendido), ya que el
                P-IGBT requiere un voltaje negativo respecto a su emisor para
                entrar en conducción.
        \end{enumerate}

  \item \textbf{Parámetros del Doble Pulso y Simulación Transitoria}
        \begin{enumerate}
          \item Configure la fuente \texttt{PULSE} para generar dos pulsos consecutivos:
                un primer pulso largo (para cargar el inductor hasta la
                corriente nominal de prueba $I_C$) y un segundo pulso corto
                (para evaluar la dinámica de conmutación) como se muestra en
                la figura~\ref{fig:pulsos}. Para el NIGBT, este debe estar
                entre 0 y 15\,V; para el PIGBT, entre 0 y $-15$\,V.
          \item Añada la directiva de simulación transitoria con una
                resolución temporal fina para capturar correctamente los
                eventos de alta frecuencia.
        \end{enumerate}

  \item \textbf{Análisis de resultados}
        \begin{enumerate}
        \item Ejecute la simulación. Grafique de forma clara $V_{GE}$,
                $V_{CE}$ y la corriente de colector $I_C$ para ambos casos.
        \item A partir del acercamiento realizado al segundo pulso de
                apagado, determine el valor de la corriente de colector
                ($I_C$) en el punto exacto donde inicia la fase asintótica
                denominada \emph{tail current}. Explique, basándose en la
                física del dispositivo, por qué esta corriente no cae
                instantáneamente a cero.
        \item Utilizando los cursores de LTspice, obtenga el valor del
                tiempo de retardo de apagado ($t_{d(off)}$) midiendo el
                intervalo entre el 90\,\% de $V_{GE}$ y el 90\,\% de $I_C$.
                Compare este dato con el valor proporcionado en la hoja de
                datos (\emph{datasheet}) del fabricante y justifique
                cualquier discrepancia encontrada.
        \item Calcule la energía de apagado ($E_{off}$) mediante la integral
                del producto $V_{CE}\cdot I_C$ durante el periodo de
                conmutación. ¿Qué porcentaje de esta energía total es
                atribuible exclusivamente al periodo de la corriente de cola?

        \item Observe el comportamiento del voltaje colector-emisor
                ($V_{CE}$) durante el apagado e identifique el valor del pico de
                sobretensión causado por las inductancias parásitas. ¿Cómo varía este pico
                de voltaje al comparar los modelos con y sin diodo de marcha libre integrado?

        \item Analice la relación entre la resistencia de compuerta
                ($R_G$) utilizada y el tiempo de retardo medido. Según sus resultados, ¿qué
                impacto tendría duplicar el valor de $R_G$ sobre las pérdidas por
                conmutación y la velocidad de respuesta del dispositivo?
        \end{enumerate}
\end{enumerate}

\begin{figure}[ht]
  \centering
  \imagen[0.5\linewidth]{fig2_esquema_igbt.png}
  \caption{Esquema a usar para los circuitos IGBT.}
  \label{fig:esquema}
\end{figure}

\begin{figure}[ht]
  \centering
  \imagen[0.5\linewidth]{fig3_pulsos_compuerta.png}
  \caption{Pulsos de control de la compuerta.}
  \label{fig:pulsos}
\end{figure}

%--------------------------------------------------------------
\section{Conclusiones}
Redactar un análisis técnico (una página como mínimo) debatiendo las
diferencias fundamentales observadas entre los canales N y P.
Posteriormente, justificar con fundamentos físicos cómo el incremento del
área del silicio (50\,A vs 75\,A) y la presencia del diodo co-empaquetado
impactan directamente en las pérdidas estáticas y dinámicas de los modelos
evaluados.

\end{document}